% concatenated from main.tex
\documentclass[12pt]{amsart}
\usepackage[utf8]{inputenc}
\usepackage{amsfonts}
\usepackage{amssymb}
\usepackage{amsmath}
\usepackage[english]{babel}
\usepackage{amsthm}
\theoremstyle{plain}
\newtheorem{theorem}{Theorem}[section]
\newtheorem{corollary}[theorem]{Corollary}
\newtheorem{lemma}[theorem]{Lemma}
\newtheorem{proposition}[theorem]{Proposition}
\theoremstyle{definition}
\newtheorem{definition}[theorem]{Definition}
\newtheorem{remark}[theorem]{Remark}
\usepackage{geometry}
\geometry{a4paper,scale=0.8}
\numberwithin{equation}{section} 
\usepackage{tikz}
\usepackage{caption}
\usepackage{graphicx}
\usepackage[titletoc,toc,title]{appendix}
\usepackage{lipsum}

\newcommand\blfootnote[1]{%
  \begingroup
  \renewcommand\thefootnote{}\footnote{#1}%
  \addtocounter{footnote}{-1}%
  \endgroup
}

\usepackage[colorlinks, citecolor=blue,pagebackref,hypertexnames=false]{hyperref}



\newcommand{\pink}{\textcolor{Magenta}}     \newcommand{\red}{\textcolor{red}}     
\newcommand{\green}{\textcolor{green}}      \newcommand{\blue}{\textcolor{blue}}    \newcommand{\yell}{\textcolor{yellow}}
\newcommand{\oran}{\textcolor{orange}}		\newcommand{\gray}{\textcolor{gray}}


%%%%%  COMMENTS, STRIKE OUT
\def\ccc#1{{\sout{#1}}}
\newcounter{comcount}

\def\comm{Comment }
%\def\comm{Komentarz }

\def\com#1{\stepcounter{comcount}
{\noindent \fboxrule=0.3mm\fcolorbox[named]{red}{white}
{\begin{minipage}{0.97\textwidth}{%\color[named]{Red}
{\bf \comm \arabic{comcount}:} #1}\end{minipage}}
\vskip 1em}}




\begin{document}
\title[Isoperimetric Inequality on Grushin space]{Isoperimetric Inequality for degenerate elliptic operators of Grushin type} 
\date{}
\author{Dangyang He}  %\sffamily
\address{Department of Mathematics, Sun Yat-sen University, Guangzhou 510275, Guangdong, P.R. China}
\email{hedy28@mail.sysu.edu.cn}






% abstract
%{\noindent\small{\bf Abstract:}
\begin{abstract}
Let \(n,m\ge 1\), \(\alpha\in(0,1)\), and \(\beta\ge 0\). For the Grushin-type operator
\[
L=-\nabla_x\!\cdot\!\bigl(|x|^{2\alpha}\nabla_x\bigr)+|x|^{2\beta}\Delta_y
\qquad \text{on } \mathbb R^n\times \mathbb R^m,
\]
we prove the isoperimetric inequality on the associated Grushin space. Equivalently, if
\[
Q=\frac{n+m(\beta+1-\alpha)}{1-\alpha},
\]
then
\[
|\Omega|^{\frac{Q-1}{Q}}\le C\,P(\Omega)
\]
for every smooth bounded domain \(\Omega\subset \mathbb R^{n+m}\).
\end{abstract}
%\vspace{1ex}

\maketitle

% main body
\tableofcontents


\blfootnote{$\textit{2020 Mathematics Subject Classification.}$ Primary 35J70; Secondary 53C17.}

\blfootnote{$\textit{Keywords and Phrases.}$ Isoperimetric inequality, Sobolev inequality, degenerate operator.}






\section{Introduction}

Let \(n,m\ge 1\), \(\alpha\in(0,1)\), and \(\beta\ge 0\). Consider the degenerate elliptic operator
\begin{equation}\label{eq:L}
L=-\nabla_x\!\cdot\!\bigl(|x|^{2\alpha}\nabla_x\bigr)+|x|^{2\beta}\Delta_y
\qquad \text{on } \mathbb R^n\times \mathbb R^m,
\end{equation}
where \(\Delta_y\) denotes the standard Laplacian in the \(y\)-variables. We refer to \cite{RS} for the precise definition and basic properties of \(L\). Associated with \(L\) is the first-order differential operator
\[
\nabla_L:=\bigl(|x|^\alpha \nabla_x,\; |x|^\beta \nabla_y\bigr),
\]
and the corresponding control distance \(d\). We write \(B(\xi,r)\) for the \(d\)-ball centered at \(\xi\) with radius \(r\), and \(P(\Omega)\) for the perimeter of a smooth set \(\Omega\subset \mathbb R^{n+m}\) with respect to this geometry.

% Assume that $n,m\ge 1$ and that $\alpha \in (0,1)$ and $\beta \ge 0$. We define a positive self-adjoint operator on $L^2(\mathbb{R}^{n+m})$ by the formula
% \begin{equation}\label{eq_Grushinoperator}
%     L = - \nabla_x \left( |x|^{2\alpha} \nabla_x \right) + |x|^{2\beta} \Delta_y,
% \end{equation}
% where $x\in \mathbb{R}^n$ and $y\in \mathbb{R}^m$. We refer readers to \cite{RS} for the precise definition and basic properties of $L$. In particular, $L$ induces a sub-Riemannian geometry. Define gradient operator $\nabla_L := \big( |x|^\alpha \nabla_x, |x|^\beta \nabla_y \big)$ and introduce Riemannian quasi-distance
% \begin{equation*}
%     d(\xi, \eta):= \sup_{\psi \in S} |\psi(\xi) - \psi(\eta)|,
% \end{equation*}
% where $S = \{\psi \in W^{1,\infty}; |\nabla_L \psi|\le 1\}$.

The purpose of this paper is to establish the isoperimetric inequality on Grushin spaces. More precisely, we prove that for every smooth bounded domain \(\Omega\subset \mathbb{R}^{n+m}\),
\begin{equation}\label{ISO_Q}
    |\Omega|^{\frac{\mathcal{Q}-1}{\mathcal{Q}}} \le C_1 P(\Omega),
\end{equation}
where \(|\cdot|\) denotes the Lebesgue measure on \(\mathbb{R}^{n+m}\), \(P(\Omega)\) denotes the perimeter of \(\Omega\), and
\begin{equation}\label{Q}
    \mathcal{Q} = \frac{n+m(\beta+1-\alpha)}{1-\alpha}
\end{equation}
is the homogeneous dimension associated with \(L\). By the co-area formula, \eqref{ISO_Q} is equivalent to the Sobolev inequality
\begin{equation}\label{S_p}\tag{$\mathrm{S}_{\mathcal{Q}}^p$}
\left(\int_{\mathbb{R}^{n+m}} |f(\xi)|^{\frac{\mathcal{Q}p}{\mathcal{Q}-p}} \, d\xi \right)^{\frac{\mathcal{Q}-p}{\mathcal{Q}}}
\le
C_p \int_{\mathbb{R}^{n+m}} |\nabla_L f(\xi)|^p \, d\xi,
\qquad \forall f\in C_c^\infty(\mathbb{R}^{n+m}),
\end{equation}
in the case \(p=1\).

% Here we present the main result of this paper

% \begin{theorem}\label{thm_ISO}
% Suppose $n, m\ge 1$, $\alpha \in (0,1)$ and $\beta \ge 0$. Then the isoperimetric inequality \eqref{ISO_Q} or \eqref{S_p} (with $p=1$) holds on $\mathbb{R}^{n+m}$. 
% \end{theorem}


% In fact, since $\nabla_L$ is local, the family \eqref{S_p} enjoys a self-improvement property
% \begin{equation*}
%     \eqref{S_p} \implies (\textrm{S}_{\mathcal{Q}}^q), \quad \forall 1\le p \le q < \mathcal{Q}.
% \end{equation*}
% It follows that

% \begin{theorem}\label{thm_S_p}
% Suppose $n, m\ge 1$, $\alpha \in (0,1)$ and $\beta \ge 0$. Then the Sobolev inequality $(\textrm{S}_{\mathcal{Q}}^p)$ holds on $\mathbb{R}^{n+m}$ for all $1\le p < \mathcal{Q}$.
% \end{theorem}

% In the Grushin plane, the isoperimetric inequality was established in \cite{MM}. In higher dimensions, Franceschi and Monti \cite{FM} proved it for $n=1$, $\alpha=0$, and $m\ge 1$; in particular, for $n\ge 2$ they obtained the result for $x$-spherically symmetric sets. Moreover, under the same assumptions as in Theorem \ref{thm_ISO}, Hauer and Sikora \cite[Theorem 6.3]{HS} proved the Sobolev inequality for $2\le p<Q$. 


We now state the main result of this paper.

\begin{theorem}\label{thm_ISO}
Suppose that \(n,m\ge 1\), \(\alpha\in(0,1)\), and \(\beta\ge 0\). Then the isoperimetric inequality \eqref{ISO_Q}, or equivalently the Sobolev inequality \eqref{S_p} with \(p=1\), holds on \(\mathbb{R}^{n+m}\).
\end{theorem}

Since \(\nabla_L\) is local, the family \eqref{S_p} enjoys a self-improvement property:
\[
(\mathrm{S}_{\mathcal Q}^p)\implies (\mathrm{S}_{\mathcal Q}^q),
\qquad 1\le p\le q<\mathcal Q.
\]
As a consequence, we obtain the following corollary.

\begin{corollary}\label{thm_S_p}
Suppose that \(n,m\ge 1\), \(\alpha\in(0,1)\), and \(\beta\ge 0\). Then the Sobolev inequality \((\mathrm{S}_{\mathcal Q}^p)\) holds on \(\mathbb{R}^{n+m}\) for all \(1\le p<\mathcal Q\).
\end{corollary}

In the Grushin plane, the isoperimetric inequality was established in \cite{MM}. In higher dimensions, Franceschi and Monti \cite{FM} proved it for \(n=1\), \(\alpha=0\), and \(m\ge 1\). Moreover, for \(n\ge 2\), they obtained the corresponding result for \(x\)-spherically symmetric sets. Finally, under the same assumptions as in Theorem~\ref{thm_ISO}, Hauer and Sikora \cite[Theorem 6.3]{HS} established the Sobolev inequality \((\mathrm{S}_{\mathcal Q}^p)\) for \(2\le p<\mathcal Q\).

% Our method is motivated by the following observation. Consider $\mathbb{R}^{n+m}$ equipped with the metric $g = |x|^{-2\alpha}dx^2 + |x|^{-2\beta}dy^2$. Transfer to polar coordinates and change variable $\rho = (1-\alpha)^{-1}r^{1-\alpha}$. It can be viewed as the doubly warped product space $(0,\infty)_\rho \times_{(1-\alpha)\rho} \mathbb{S}^{n-1} \times_{[(1-\alpha)\rho]^{-\beta/(1-\alpha)}} \mathbb{R}^m$ with metric 
% $$g = d\rho^2 + [(1-\alpha)\rho]^{2} d\theta^2 + [(1-\alpha)\rho]^{-2\beta/(1-\alpha)}dy^2.$$
% By the standard formula for doubly warped product spaces, one can estimate the Ricci curvature of $g$:
% \begin{equation*}
%     \textrm{Ric}_\xi \ge - \frac{C(n,m,\alpha,\beta)}{\rho^2} g_\xi = - \frac{C(n,m,\alpha,\beta)}{|x|^{2-2\alpha}} g_\xi.
% \end{equation*}
% Note that our Grushin operator $L$ is not the standard Laplace--Beltrami operator associated with $g_\xi$. Nevertheless, the above calculation suggests a corresponding curvature-dimension inequality. Combined with a remote ball argument from \cite{Gilles}, this leads to a suitable gradient estimate for the heat semigroup. We then apply the method of \cite{He_ISO}, based on a weak-type Sobolev inequality and a Hardy-type inequality, to obtain the desired isoperimetric estimate.


Our argument is motivated by the following observation. Endow \(\mathbb{R}^{n+m}\) with the metric
\[
g=|x|^{-2\alpha}dx^2+|x|^{-2\beta}dy^2.
\]
Passing to polar coordinates in the \(x\)-variable and introducing $\rho=(1-\alpha)^{-1}r^{1-\alpha}$, one sees that \(g\) can be written as the doubly warped product metric
\[
g=d\rho^2+[(1-\alpha)\rho]^2\,d\theta^2+[(1-\alpha)\rho]^{-2\beta/(1-\alpha)}\,dy^2
\]
on
\[
(0,\infty)_\rho\times_{(1-\alpha)\rho}\mathbb{S}^{n-1}\times_{[(1-\alpha)\rho]^{-\beta/(1-\alpha)}}\mathbb{R}^m.
\]
By the standard curvature formula for doubly warped products,
\[
\textrm{Ric}_\xi \ge -\frac{C(n,m,\alpha,\beta)}{\rho^2}\,g_\xi
= -\frac{C(n,m,\alpha,\beta)}{|x|^{2-2\alpha}}\,g_\xi.
\]
Although our Grushin operator \(L\) is not the standard Laplace--Beltrami operator associated with \(g\), this computation suggests a corresponding curvature-dimension inequality for \(-L\). Combined with a remote ball argument from \cite{Gilles}, this yields a suitable gradient estimate for the heat semigroup. We then apply the method of \cite{He_ISO}, based on a weak-type Sobolev inequality and a Hardy-type inequality, to obtain the desired isoperimetric estimate.








\section{Preliminaries}

Throughout the paper, we write \(\xi=(x,y)\) and \(\eta=(x',y')\), where \(x,x'\in\mathbb R^n\) and \(y,y'\in\mathbb R^m\). We denote by \(B(\xi,r)\) the Grushin ball $B(\xi,r):=\{\eta\in\mathbb R^{n+m}: d(\xi,\eta)<r\}$, and by $V(\xi,r):=|B(\xi,r)|$ its volume with respect to the Lebesgue measure on \(\mathbb R^{n+m}\). We also write \(B_n\) and \(B_m\) for Euclidean balls in \(\mathbb R^n\) and \(\mathbb R^m\), respectively.

The following lemma summarizes some basic geometric properties of Grushin spaces; see \cite[Proposition~5.1, Theorem~6.4]{RS}.


\begin{lemma}\label{le_geo}
Let $n, m\ge 1$, $\alpha \in (0,1)$ and $\beta \ge 0$.

\begin{enumerate}
\item Distance estimate: there exists constant $C_d>0$ such that \begin{align}
   C_d^{-1} D(\xi,\eta) \le  d(\xi,\eta) \le C_d D(\xi,\eta),
\end{align}
where
\begin{equation*}
    D(\xi,\eta) = \frac{|x-x'|}{\left( |x| + |x'| \right)^\alpha} + \frac{|y-y'|}{(|x|+|x'|)^{\beta} + |y-y'|^{\frac{\beta}{\beta+1-\alpha}}}.
\end{equation*}
\item Volume estimate: \begin{align*}
    V(\xi,r) = V((x,y),r) \sim \begin{cases}
        r^{\mathcal{Q}}, & r\ge |x|^{1-\alpha},\\
        r^{n+m}|x|^{n \alpha + m\beta}, & r\le |x|^{1-\alpha}.
    \end{cases}
\end{align*}
In addition, there exists constant $C_D>0$ such that the volume doubling condition holds:
\begin{align*}
\frac{V(\xi, R)}{V(\xi,r)} \le C_D \left(\frac{R}{r}\right)^{\mathcal{Q}}
\end{align*}
for all $\xi \in \mathbb{R}^{n+m}$ and all $R\ge r>0$.
\item Gaussian upper bound for heat kernel: 
\begin{equation*}
    e^{-tL}(\xi,\eta) \le \frac{C}{V(\xi, \sqrt{t})} e^{-\frac{d(\xi,\eta)^2}{ct}}
\end{equation*}
for all $t>0$ and almost every $\xi,\eta \in \mathbb{R}^{n+m}$.
\end{enumerate}

\end{lemma}

With Lemma~\ref{le_geo}, we develop the following property.


\begin{lemma}\label{le_volume}
Let $n,m\ge 1$, $\alpha\in (0,1)$ and $\beta\ge 0$. Let $\xi=(x,y)\in \mathbb{R}^{n+m}$. There exist $0<c_1<c_2$ and $\epsilon_0 \in (0,1)$ such that for all $r\le \epsilon_0 |x|^{1-\alpha}$, one has
\begin{equation}
    B_n(x, c_1 r |x|^\alpha) \times B_m(y, c_1 r |x|^\beta) \subset B(\xi, r) \subset B_n(x, c_2 r |x|^\alpha) \times B_m(y, c_2 r |x|^\beta).
\end{equation}
Moreover, there exist $0<c_3<c_4$ such that for all $r>0$, one has for all $y\in \mathbb{R}^m$
\begin{equation}\label{le_volume3}
    B_n(0, c_3 r) \times B_m(y, c_3 r^{\beta+1-\alpha}) \subset B((0,y), r^{1-\alpha}) \subset B_n(0, c_4 r) \times B_m(y, c_4 r^{\beta+1-\alpha}).
\end{equation}
\end{lemma}


\begin{proof}[Proof of Lemma~\ref{le_volume}]
We first show that 
\begin{equation}\label{le_volume1}
    B(\xi, r) \subset B_n(x, c_2 r |x|^\alpha) \times B_m(y, c_2 r |x|^\beta).
\end{equation}
Take $\eta = (x',y')\in B(\xi, r)$. Then by Lemma~\ref{le_geo}, $D(\xi,\eta)\le C_d r$. We claim that
\begin{equation}\label{claim1}
    |x-x'| \le \frac{|x|}{2},
\end{equation}
provided $\epsilon_0$ is small enough. Indeed, if not, then $|x|+|x'|\le |x-x'| + 2|x| \le 3|x-x'|$ and thus
\begin{equation*}
    \frac{|x-x'|}{\left(|x|+|x'|\right)^\alpha} \ge 3^{-\alpha} |x-x'|^{1-\alpha} \ge 3^{-\alpha} 2^{\alpha-1} |x|^{1-\alpha},
\end{equation*}
which further implies that $D(\xi,\eta)\ge 3^{-\alpha} 2^{\alpha-1} |x|^{1-\alpha}$. Pick $\epsilon_0 < C_d^{-1} 3^{-\alpha} 2^{\alpha-1}$. Then the previous argument contradicts the fact that $D(\xi,\eta)\le C_d r$ and hence the claim \eqref{claim1} is verified. With \eqref{claim1}, one gets
\begin{equation}\label{eq_volume2}
    \frac{|x|}{2} \le |x'| \le \frac{3}{2}|x|\quad \textrm{and}\quad |x|+|x'| \le \frac{5}{2} |x|.
\end{equation}

Next, since $D(\xi,\eta)\le C_d r$, we must have
\begin{align*}
    \frac{|x-x'|}{\left(|x|+|x'|\right)^\alpha} \le C_d r \implies |x-x'|\le C_d r \left(|x|+|x'|\right)^\alpha \le C_d \left(\frac{5}{2}\right)^\alpha r |x|^\alpha.
\end{align*}
While for the $y$-part of $D(\xi,\eta)$, if $\left(|x|+|x'|\right)^{\beta+1-\alpha}\ge |y-y'|$ then one obtains
\begin{align*}
\frac{|y-y'|}{2\left(|x|+|x'|\right)^{\beta}} \le \frac{|y-y'|}{\left(|x|+|x'|\right)^{\beta} + |y-y'|^{\frac{\beta}{\beta+1-\alpha}}} \le C_d r.
\end{align*}
By \eqref{eq_volume2}, one yields
\begin{equation*}
    |y-y'|\le 2C_d r \left(|x|+|x'|\right)^{\beta} \le 2C_d \left(\frac{5}{2}\right)^\beta r |x|^\beta.
\end{equation*}
If $\left(|x|+|x'|\right)^{\beta+1-\alpha}\le |y-y'|$, then
\begin{align*}
|y-y'|^{\frac{1-\alpha}{\beta+1-\alpha}} \le 2 C_d r \implies |y-y'|\le \left(2C_d\right)^{\frac{\beta+1-\alpha}{1-\alpha}} r r^{\frac{\beta}{1-\alpha}}\le \left(2C_d\right)^{\frac{\beta+1-\alpha}{1-\alpha}} r |x|^{\beta}.
\end{align*}
Setting $c_2 = \max \left( C_d \left(\frac{5}{2}\right)^\alpha,  2C_d \left(\frac{5}{2}\right)^\beta, \left(2C_d\right)^{\frac{\beta+1-\alpha}{1-\alpha}} \right)$, one concludes \eqref{le_volume1}.

The proof of the converse direction:
\begin{equation}
    B_n(x, c_1 r |x|^\alpha) \times B_m(y, c_1 r |x|^\beta) \subset B(\xi, r)
\end{equation}
is obvious. Indeed, for $ \eta = (x',y')$ such that $|x-x'|\le c_1 r |x|^\alpha$ and $|y-y'|\le c_1 r |x|^\beta$, one directly bounds
\begin{equation*}
    d(\xi,\eta)\le C_d D(\xi,\eta)\le C_d \frac{c_1 r |x|^\alpha}{\left(|x|+|x'|\right)^\alpha} + C_d \frac{c_1 r |x|^\beta}{\left(|x|+|x'|\right)^\beta+|y-y'|^{\frac{\beta}{\beta+1-\alpha}}} \le 2 C_d c_1 r.
\end{equation*}
Choosing $c_1 = \frac{1}{4C_d}$, one confirms $\eta \in B(\xi,r)$ as desired.

The proof for \eqref{le_volume3} is similar and we omit it.
\end{proof}




\begin{remark}
Note that the constants \(c_1\) and \(c_2\) depend only on \(C_d\), \(\alpha\), and \(\beta\). Therefore, by choosing $\epsilon_0< \min \left( C_d^{-1} 3^{-\alpha} 2^{\alpha-1}, \frac{1}{2c_2} \right)$, one sees that for every \(r\le \epsilon_0 |x|^{1-\alpha}\) and every $\eta=(x',y')\in B_n(x, c_2 r |x|^\alpha)\times B_m(y, c_2 r |x|^\beta)$, we have
\[
|x|\le |x-x'|+|x'|
\le c_2 r |x|^\alpha + |x'|
\le c_2\epsilon_0 |x|+|x'|
\le \frac{|x|}{2}+|x'|.
\]
Hence $|x|\le 2|x'|\le 3|x|$.
\end{remark}
    

\begin{definition}
A ball \(B=B(\xi,r)=B((x,y),r)\) is said to be \emph{remote} if
\[
r\le \epsilon_0 |x|^{1-\alpha},
\]
where \(\epsilon_0\) is given by Lemma~\ref{le_volume}. For such a ball, every \(\eta=(x',y')\in B\) satisfies
\[
|x-x'|\le \frac{|x|}{2}
\qquad\text{and hence}\qquad
|x|\sim |x'|.
\]
\end{definition}



Next, we recall the following result from \cite[Theorem~3.1, Proposition~3.5]{BCLS}.


 
\begin{lemma}\label{le_BCLS}
Let $0<p<\infty$, $r_0,s_0\in (0,\infty]$ and $\theta_0 \in (0,1]$ and that the parameter $q = q(r_0,s_0,\theta_0)$ defined by
\begin{align*}
    \frac{1}{r_0}  = \frac{\theta_0}{q} + \frac{1-\theta_0}{s_0},
\end{align*}
satisfies $p\le q<\infty$. Suppose that
\begin{equation}
    \sup_{\lambda>0} \lambda  \left| \{|f|\ge \lambda\}\right|^{\frac{1}{r_0}}  \le C \|\nabla_L f\|_p^{\theta_0} \|f\|_\infty^{1-\theta_0} |\textrm{supp}(f)|^{\frac{1-\theta_0}{s_0}}.
\end{equation}
Then, we have
\begin{align*}
    \|f\|_q \le C \| \nabla_L f\|_p,\quad q = \frac{r_0 s_0 \theta_0}{s_0 - r_0(1-\theta_0)}
\end{align*}
for all $f\in C_c^\infty(\mathbb{R}^{n+m})$.
\end{lemma}

\begin{remark}
In particular, if we take $p=1$ and show that
\begin{equation}\label{showthis}
\sup_{\lambda>0} \lambda^{\frac{\mathcal{Q}}{\mathcal{Q}-1}}  \left|\left\{ |f| \ge \lambda \right\}\right| \le C \|\nabla_L f\|_1 \|f\|_q^{\frac{1}{\mathcal{Q}-1}},
\end{equation}
then Lemma~\ref{le_BCLS} yields $(\textrm{S}_{\mathcal{Q}}^1)$ and hence \eqref{ISO_Q}.
\end{remark}




\section{Proof of Theorem~\ref{thm_ISO}}

Given a symmetric Markov generator $\mathcal{L}$ with carré du champ $\Gamma$ and iterate $\Gamma_2$, the inequality
\begin{align*}
    \Gamma_2(f) \ge K \Gamma(f) + \frac{1}{N} \left(L(f)\right)^2,\quad \forall f\in C^\infty
\end{align*}
encodes both a 'Ricci curvature' lower bound $K$ and 'effective dimension' $N$. For conventional issues, we let $\mathcal{L}$ be a \textbf{non-positive} self-adjoint diffusion operator. Formally, for smooth function $f,g$ one defines bilinear form:
\begin{align}\label{carre1}
    \Gamma(f,g) = \frac{1}{2} \left( \mathcal{L}(fg) - f \mathcal{L}(g) - g\mathcal{L}(f)\right).
\end{align}
Note that $\Gamma$ is symmetric and $\Gamma(f) = \Gamma(f,f)$ is the so-called \textit{carre du champ} operator. Next, one defines the iterated bilinear form:
\begin{align}\label{carre2}
    \Gamma_2(f,g) = \frac{1}{2} \left( \mathcal{L}\Gamma(f,g) - \Gamma(f,\mathcal{L}(g)) - \Gamma(g,\mathcal{L}(f)) \right),
\end{align}
and $\Gamma_2(f) = \Gamma_2(f,f)$ is the second \textit{carre du champ} operator.





\begin{lemma}\label{le_Ric_Grushin}
Let $n,m\ge 1$, $\alpha \in (0,1)$ and $\beta \ge 0$. Let $\Gamma$ (resp. $\Gamma_2$) be the \textit{carre du champ} (resp. second \textit{carre du champ}) of $L$. The following curvature dimension inequality holds
\begin{align*}
\Gamma_2(f) \ge - \frac{C(n,m,\alpha,\beta)}{|x|^{2(1-\alpha)}} \Gamma(f) + C'(n,m,\alpha,\beta) \left(L(f) \right)^2
\end{align*}
for all $f\in C^\infty( \mathbb{R}^{n}\setminus\{0\} \times \mathbb{R}^m)$.
\end{lemma}

\begin{proof}[Proof of Lemma~\ref{le_Ric_Grushin}]
Since $L$ is non-negative, we set $\mathcal{L} = -L$ and calculate \eqref{carre1} and \eqref{carre2} directly. Put $r=|x|$, $a = r^{2\alpha}$ and $b=r^{2\beta}$. Then
\begin{equation*}
\mathcal L = -a\Delta_x - b\Delta_y+\nabla a\cdot \nabla_x.
\end{equation*}

For simplicity, we use indices $i,j\in\{1,\dots,n\}$, $\mu,\nu\in\{1,\dots,m\}$, one gets directly from the definition \eqref{carre1} that
\begin{equation}\label{eq_gamma1}
    \Gamma(f)=a f_i^2+b f_\mu^2.
\end{equation}

Next, by \eqref{carre2}, a direct calculation gives
\begin{align}\label{eq_gamma2}
\Gamma_2(f)
={}&a^2 f_{ij}^2+2ab f_{i\mu}^2+b^2 f_{\mu\nu}^2\\ \nonumber
&+2a a_i f_j f_{ij} + a a_j f_j \Delta_x f\\ \nonumber
&+2a b_i f_\mu f_{i\mu}+ a b_i f_i \Delta_y f\\ \nonumber
&+\Bigl(-\frac a2 \Delta_x a+\frac12|\nabla a|^2\Bigr)f_i^2
-a a_{ij}f_i f_j\\ \nonumber
&+\Bigl(-\frac a2 \Delta_x b+\frac12\nabla a\cdot \nabla b\Bigr)f_\mu^2.
\end{align}
Note that
\begin{align*}
    a_i=2\alpha r^{2\alpha-2}x_i,\quad
|\nabla a|^2 = 4\alpha^2 r^{4\alpha-2},
\quad
\Delta_x a = -2\alpha(n+2\alpha-2)r^{2\alpha-2},
\end{align*}
\begin{equation*}
    a_{ij} = 2\alpha r^{2\alpha-2}\delta_{ij}
 + 4\alpha(\alpha-1)r^{2\alpha-4}x_i x_j,
\end{equation*}
and 
\begin{align*}
    b_i = 2\beta r^{2\beta-2}x_i,\quad
|\nabla b|^2 = 4\beta^2 r^{4\beta-2},
\quad
\Delta_x b = -2\beta(n+2\beta-2)r^{2\beta-2}.
\end{align*}


We estimate the lower-order terms in \eqref{eq_gamma2} by Young's inequality.

First, using $|\Delta_x f|\le \sqrt n|D_x^2 f|$, where $|D_x^2f|^2 = \sum_{1\le i,j\le n}|\partial_{x_i,x_j}f|^2$, and Cauchy-Schwartz inequality and then Young's inequality, one obtains
\begin{equation}\label{1}
    2a a_i f_j f_{ij} + a a_j f_j \Delta_x f \ge
-\frac14 a^2|D_x^2 f|^2-C|\nabla a|^2|\nabla_x f|^2.
\end{equation}
Also, since $\alpha<1$,
\begin{equation}\label{2}
-a a_{ij}f_i f_j = -2\alpha r^{4\alpha-2}|\nabla_x f|^2 + 4\alpha(1-\alpha)r^{4\alpha-4}(x\cdot \nabla_x f)^2
\ge -2\alpha r^{4\alpha-2}|\nabla_x f|^2.
\end{equation}
Hence
\begin{equation}\label{3}
    \Bigl(-\frac a2 \Delta_x a + \frac12| \nabla a|^2\Bigr)f_i^2-a a_{ij}f_i f_j \ge
-C r^{4\alpha-2}|\nabla_x f|^2.
\end{equation}

Next, it is clear that
\begin{equation}\label{4}
    2a b_i f_\mu f_{i\mu}
\ge
-a b f_{i\mu}^2-\frac{a|\nabla b|^2}{b} |\nabla_y f|^2,
\end{equation}
and, using $|\Delta_y f|\le \sqrt m|D_y^2 f|$, where $|D_y^2f|^2 = \sum_{1\le \mu,\nu\le m}|\partial_{y_\mu,y_\nu}f|^2$, 
\begin{equation}\label{5}
    a b_i f_i \Delta_y f \ge -\frac14 b^2|D_y^2 f|^2 - C \frac{a^2|\nabla b|^2}{b^2}|\nabla_x f|^2.
\end{equation}

Finally,
\begin{equation}\label{6}
    \Bigl( -\frac a2 \Delta_x b + \frac12 \nabla a \cdot \nabla b \Bigr) f_\mu^2 \ge
 -  C r^{2\alpha+2\beta-2}|\nabla_y f|^2.
\end{equation}

Insert \eqref{1}--\eqref{6} into \eqref{eq_gamma2}. 
% \begin{equation*}
%     |\nabla a|^2\sim r^{4\alpha-2},\quad
% \frac{a|\nabla b|^2}{b}\sim r^{2\alpha+2\beta-2},\quad
% \frac{a^2|\nabla b|^2}{b^2}\sim r^{4\alpha-2},
% \end{equation*}
We obtain
\begin{align}\label{7}
{}&\Gamma_2(f)\\ \nonumber
&\ge \frac34 a^2|D_x^2 f|^2 + ab|D_{xy}^2 f|^2 + \frac34 b^2|D_y^2 f|^2 - C \Bigl(r^{4\alpha-2}|\nabla_x f|^2 + r^{2\alpha+2\beta-2}|\nabla_y f|^2\Bigr)\\ \nonumber
&= \frac34 a^2|D_x^2 f|^2 + ab|D_{xy}^2 f|^2 + \frac34 b^2|D_y^2 f|^2 -\frac{C}{|x|^{2-2\alpha}}\Gamma(f)
\end{align}

On the other hand, by trace inequality, with $C'=12(n+m)$, 
\begin{align}\label{8}
C'^{-1}(\mathcal L f)^2 &\le
C'^{-1} 3a^2(\Delta_x f)^2+ C'^{-1} 3b^2(\Delta_y f)^2+ C'^{-1}3|\nabla a|^2|\nabla_x f|^2\\ \nonumber
&\le C'^{-1} 3n a^2|D_x^2 f|^2 + C'^{-1} 3m b^2|D_y^2 f|^2 + C'^{-1} 12\alpha^2 r^{4\alpha-2}|\nabla_x f|^2\\ \nonumber
&\le \frac14 a^2|D_x^2 f|^2+\frac14 b^2|D_y^2 f|^2+\frac1{|x|^{2-2\alpha}}\Gamma(f).
\end{align}

Combining \eqref{7} and \eqref{8}, we conclude that
\begin{equation*}
    \Gamma_2(f)\ge -\frac{C}{|x|^{2-2\alpha}}\Gamma(f)+\frac1{C'}(\mathcal L f)^2
\end{equation*}
for some $C=C(n,m,\alpha,\beta)>0$, as desired.





\end{proof}




Next, we establish a gradient estimate for the heat semigroup. The argument is inspired by \cite[Sections~3.2--3.3]{Gilles}.




\begin{lemma}\label{le_Grushin_gradient}
Let $n, m\ge 1$, $\alpha \in (0,1)$ and $\beta \ge 0$. The following gradient estimate holds
\begin{align*}
    \left|\nabla_{L}e^{-tL}(\xi,\eta)\right|\le \left(\frac{1}{\sqrt{t}} + \frac{1}{|x|^{1-\alpha}}\right) \frac{C}{V(\xi,\sqrt{t})}e^{-\frac{d(\xi,\eta)^2}{ct}},
\end{align*}
where $\xi = (x,y) \in \mathbb{R}^n \setminus \{0\} \times \mathbb{R}^m$.
\end{lemma}

% \begin{proof}[Proof of Lemma~\ref{le_Grushin_gradient}]
% Let $\xi = (x,y) \in \mathbb{R}^n \setminus \{0\} \times \mathbb{R}^m$. Let $B = B(\xi,r)$ be a remote ball with radius $r \le \epsilon_0|x|^{1-\alpha}$. By Lemma~\ref{le_Ric_Grushin}, one checks easily that for any $\xi' \in B$, the curvature dimension inequality, $\textrm{CD}(-c_1/r^2, c_2)$, holds uniformly for some constants $c_1,c_2$. Therefore, by \cite[Theorem~1.4]{ZZ}; see also \cite[Theorem~5.1]{XDL}, a local version Li-Yau type estimate indicates that for $\xi' \in B/2$, $\gamma >1$, and any positive $u_t$ such that $\partial_t u_t = - Lu_t$, one has
% \begin{align}\label{le_grushin_grad1}
%     \frac{|\nabla_L u_t|^2}{u_t^2} - \gamma \frac{\partial_t u_t}{u_t} \le C(\gamma,n,m) \left(\frac{1}{t} + \frac{1}{r^2} \right).
% \end{align}
% Combining this and \cite[Theorem~2.6]{Sturm3} (say), i.e., 
% \begin{align}\label{le_grushin_grad2}
%     \left| \partial_t e^{-tL}(\xi,\eta)\right| \le \frac{C}{t V(\xi,\sqrt{t})} e^{-\frac{d(\xi,\eta)^2}{ct}}
% \end{align}
% yields the gradient estimate as desired.

% \end{proof}

\begin{proof}[Proof of Lemma~\ref{le_Grushin_gradient}]
Let \(\xi=(x,y)\in \mathbb{R}^n\setminus\{0\}\times \mathbb{R}^m\), and let \(B=B(\xi,r)\) be a remote ball with radius \(r\le \epsilon_0 |x|^{1-\alpha}\). By Lemma~\ref{le_Ric_Grushin}, one readily checks that for every \(\xi'\in B\), the curvature-dimension inequality \(\mathrm{CD}(-c_1/r^2,c_2)\) holds uniformly for some constants \(c_1,c_2>0\). Therefore, by \cite[Theorem~1.4]{ZZ}; see also \cite[Theorem~5.1]{XDL}, a local Li--Yau type estimate implies that for every \(\xi'\in B/2\), every \(\gamma>1\), and every positive solution \(u_t\) of
\[
\partial_t u_t=-Lu_t,
\]
one has
\begin{equation}\label{le_grushin_grad1}
    \frac{|\nabla_L u_t|^2}{u_t^2}-\gamma \frac{\partial_t u_t}{u_t}
    \le
    C(\gamma,n,m)\left(\frac{1}{t}+\frac{1}{r^2}\right).
\end{equation}
Combining \eqref{le_grushin_grad1} with \cite[Theorem~2.6]{Sturm3}, namely
\begin{equation}\label{le_grushin_grad2}
    \left| \partial_t e^{-tL}(\xi,\eta)\right|
    \le
    \frac{C}{t\,V(\xi,\sqrt t)}e^{-\frac{d(\xi,\eta)^2}{ct}},
\end{equation}
yields the desired gradient estimate.
\end{proof}







The next step is to establish the following Hardy type inequality.
\begin{lemma}\label{le_hardy}
Let $n, m\ge 1$, $\alpha \in (0,1)$ and $\beta \ge 0$. Then, the following Hardy type inequality holds
\begin{align}\label{eq_hardy}
    \int_{\mathbb{R}^{n+m}} \frac{|u(\xi)|}{|x|^{1-\alpha}} d\xi \le C_H \int_{\mathbb{R}^{n+m}} |\nabla_{L}u(\xi)| d\xi,
\end{align}
for all $u\in C_c^\infty(\mathbb{R}^{n+m})$.
\end{lemma}


\begin{proof}[Proof of Lemma~\ref{le_hardy}]
We give a direct proof. Assume first that \(n\ge 2\). Passing to polar coordinates in the \(x\)-variable, the left-hand side of \eqref{eq_hardy} becomes
\begin{align*}
    C(n) \int_{\mathbb{R}^m} \int_0^\infty \int_{\mathbb{S}^{n-1}}
    \frac{|u((r,\theta);y)|}{r^{1-\alpha}}\, d\theta\, r^{n-1}\, dr\, dy.
\end{align*}
Using the inequality
\begin{align*}
    |u((r,\theta);y)|
    \le
    \int_r^\infty |\partial_s u((s,\theta);y)|\, ds,
\end{align*}
we deduce that the left-hand side of \eqref{eq_hardy} is bounded by
\begin{align*}
    C(n) \int_{\mathbb{R}^m} \int_{\mathbb{S}^{n-1}} \int_0^\infty \int_r^\infty
    |\partial_s u((s,\theta);y)|\, ds\, r^{n-2+\alpha}\, dr\, d\theta\, dy.
\end{align*}
Since \(r\le s\) in the region of integration, we have \(r^\alpha\le s^\alpha\). Moreover,
\[
|s^\alpha \partial_s u|\le |\nabla_L u|.
\]
Therefore, the above expression is bounded by
\begin{align*}
    C(n) \int_{\mathbb{R}^m} \int_{\mathbb{S}^{n-1}} \int_0^\infty \int_0^s
    |\nabla_L u((s,\theta);y)|\, r^{n-2}\, dr\, ds\, d\theta\, dy
    \\
    \le
    C(n) \int_{\mathbb{R}^m} \int_{\mathbb{S}^{n-1}} \int_0^\infty
    |\nabla_L u((s,\theta);y)|\, s^{n-1}\, ds\, d\theta\, dy
    \\
    \le
    C(n) \int_{\mathbb{R}^{n+m}} |\nabla_L u(\xi)|\, d\xi.
\end{align*}
The case \(n=1\) is similar, and we omit the details.
\end{proof}






% \begin{proof}[Proof of Lemma~\ref{le_hardy}]
% We give a direct proof. Suppose first $n\ge 2$. With polar coordinates, the LHS of \eqref{eq_hardy} equates
% \begin{align*}
%     C(n) \int_{\mathbb{R}^m} \int_0^\infty \int_{\mathbb{S}^{n-1}} \frac{|u((r,\theta); y)|}{r^{1-\alpha}} d\theta r^{n-1} dr dy.
% \end{align*}
% With inequality
% \begin{align*}
%     |u((r,\theta);y)|\le \int_r^\infty |\partial_s u((s,\theta);y)| ds,
% \end{align*}
% one deduces that the LHS of \eqref{eq_hardy} is bounded by
% \begin{align*}
%     C(n) \int_{\mathbb{R}^m} \int_{\mathbb{S}^{n-1}} \int_0^\infty \int_r^\infty |\partial_s u((s,\theta);y)| ds r^{n-2+\alpha} dr d\theta dy.
% \end{align*}
% Trivially, one bounds $r^\alpha \le s^\alpha$ and notice that $|s^\alpha \partial_s u|\le |\nabla_L u|$. One estimates the above by
% \begin{align*}
%     C(n) \int_{\mathbb{R}^m} \int_{\mathbb{S}^{n-1}} \int_0^\infty &\int_0^s |\nabla_L u((s,\theta);y)| r^{n-2} dr ds d\theta dy\\
%     &\le C(n) \int_{\mathbb{R}^m} \int_{\mathbb{S}^{n-1}} \int_0^\infty |\nabla_L u((s,\theta);y)| s^{n-1} ds d\theta dy\\
%     &\le C(n) \int_{\mathbb{R}^{n+m}} |\nabla_L u(\xi)| d\xi.
% \end{align*}

% The case $n=1$ is similar and we omit details. 

%Indeed, split the integral into $\int_{\mathbb{R}} = \int_0^\infty + \int_{-\infty}^0$ and estimate them separately. One checks easily that
% \begin{align*}
%     \int_{\mathbb{R}^m} \int_0^\infty \frac{|u(r,y)|}{|r|^{1-\alpha}} dr dy \le C \int_{\mathbb{R}^m} \int_0^\infty |\nabla_L u(s,y)| dsdy,
% \end{align*}
% and
% \begin{align*}
%     \int_{\mathbb{R}^m} \int_{-\infty}^0 \frac{|u(r,y)|}{|r|^{1-\alpha}} dr dy \le C \int_{\mathbb{R}^m} \int_{-\infty}^0 |\nabla_L u(s,y)| dsdy.
% \end{align*}
% The proof follows easily.
% \end{proof}




The next lemma is a variant of the argument in \cite[Section~2.1]{DR_hardy}. We include a proof for completeness.

\begin{lemma}\label{le_remoteball}
Let $R>0$. There exists a sequence of balls $\{B_n^j = B(x_j, r_j)\}_{j \ge 0}$ such that 
\begin{enumerate}
    \item $\mathbb{R}^n = \cup_{j\ge 0} B_n^j$, where $B_n^0 = B_n(0,R)$ and $B_j$ ($j\ge 1$) is $x$-remote in the sense $r_j \le |x_j|/2$,
    \item for all $j \ge 1$, $2^{-10} |x_j| \le r_j \le 2^{-9} |x_j|$,
    \item $\sum_j \mathbf{1}_{B_j}(x) \le C_L$ for all $x\in \mathbb{R}^n$, where $C_L>0$ does not depend on $R$.
\end{enumerate}

\end{lemma}

\begin{proof}[Proof of Lemma~\ref{le_remoteball}]
Set $B_n^0 = B_n(0,R)$ and $A_N:= B_n(0, R 2^{N}) \setminus B_n(0, R 2^{N-1})$ for each $N\ge 1$. Apparently, 
\begin{align*}
    \mathbb{R}^n = B_n^0 \cup \bigcup_{N\ge 1}A_N,\quad A_N \subset \bigcup_{x\in A_N} B_n\big(x, R 2^{N-13}\big).
\end{align*}
It follows by Vitali's covering lemma, that one can find an index set $I_N$ and a collection of balls $\{B_n\big(x_{N,j},R2^{N-13}\big)\}_{j\in I_N}$ with $x_{N,j}\in A_N$, pairwise disjoint and $A_N \subset \bigcup_{j\in I_N}B_n\big(x_{N,j}, R2^{N-10}\big)$. One deduces the finiteness of $\# I_N$, 
\begin{align*}
    \# I_N |B_n(0,R 2^{N+2})| &\le \sum_{j \in I_N} |B_n\big(x_{N,j}, R2^{N+3}\big)|\\
    &\le  2^{16n} \sum_{j \in I_N} |B_n(x_{N,j}, R2^{N-13})| \\
    &\le 2^{16n}|B_n \big(0,R2^{N+2}\big)|,
\end{align*}
Note that by setting $B_n^{j} = B_n(x_{j}, r_j)$ with $x_j = x_{N,j}$, $r_j = R 2^{N-10}$ and relabeling, we construct a sequence of balls $\{B_n^j\}_j$ such that $\cup_{j \ge 0}B_n^j = \mathbb{R}^n$. Moreover, we have for $j \ne 0$ (since $x_j \in A_N$ for some $N$),
\begin{align*}
2^{-10}|x_j|\le r_j \le 2^{-9} |x_j|,
\end{align*}
and for all $x\in B_n^j$,
\begin{align*}
    2^8 r_j \le |x| \le 2^{11}r_j.
\end{align*}
Clearly, this construction guarantees that $B_n^0 = B_n(0,R)$ is the only ball in the collection which contains $0$. In addition, for $x\ne 0$, one sets $J_x = \{j; x\in B_n^j\}$. Observe that if $x\in A_N$ for some $N\ge 1$ (the case $x\in B_n^0$ is similar), and $x\in B_n^j$ for some $j$, then $j$ is either in $I_N$, $I_{N-1}$ or $I_{N+1}$. Therefore, if $x\in A_N$, we have
\begin{align*}
    \# J_x &\big|B_n\big(x,2^{-1}|x| \big)\big| \le \sum_{j \in J_x} \big|B_n \big(x_j, 2^{-1}|x|+2^{-8}|x|\big)\big| \le 2^{20n} \sum_{j\in J_x} \big|B_n \big(x_j, 2^{-20}|x|\big) \big|\\
    &\le 2^{20n} \sum_{j\in J_x} \big|B_n\big(x_j, 2^{-3}r_j\big)\big| = 2^{20n} \sum_{l\in \{-1,0,1\}} \sum_{j \in J_x \cap I_{N+l}} \big|B_n\big(x_j, 2^{-3}r_j\big)\big| \\
    &=2^{20n} \sum_{l\in \{-1,0,1\}} \left|\cup_{j \in J_x\cap I_{N+l}}B_n\big(x_j, 2^{-3} r_j\big)\right|\\
    &\le 3 \cdot 2^{20n} \big|B_n\big(x, 2^{-1}|x|\big)\big|,
\end{align*}
since for all $j \in J_x\cap I_{N+l}$, one has $$B_n\big(x_j, 2^{-3} r_j\big) \subset B_n\big(x,9 \cdot 2^{-3} r_j\big) \subset B_n \big(x, 9\cdot 2^{-11} |x| \big) \subset B_n\big(x, 2^{-1}|x|\big).$$ 
The case $x\in B_n^0$ is similar. Hence, we conclude that there exists a constant $C_L>0$ which does not depend on $R,N$ such that
\begin{align*}
    \sum_j \mathbf{1}_{B_j}(x) \le C_L,\quad \forall x\in \mathbb{R}^n,
\end{align*}
completing the proof.
\end{proof}






We are now in a position to prove Theorem~\ref{thm_ISO}


\begin{proof}[Proof of Theorem~\ref{thm_ISO}]
    

Let $\lambda>0$, $q=\frac{\mathcal{Q}}{\mathcal{Q}-1}$ and $f\in C_c^\infty(\mathbb{R}^{n+m})$ be fixed. Set $R = \left(\frac{\|f\|_q}{\lambda}\right)^{\frac{q}{\mathcal{Q}(1-\alpha)}}$. By Lemma~\ref{le_remoteball} with parameter $R$, there exists a sequence of balls $\{B_n^{j}=B_n(x_j,r_j)\}_{j \ge 0}$ such that
\begin{align*}
    \mathbb{R}^n = B_n^0 \cup \left(\cup_{j \ge 1}B_n^{j} \right).
\end{align*}
Moreover, each $B_n^{j}$ ($j \ne0$) is $x$-remote and $B_n^0 = B_n(0,R)$. Let $\{\mathcal{X}_\alpha\}$ be a smooth partition of unity subordinated to $B_n^{j}$ such that $0\le \mathcal{X}_j \le 1$ and
\begin{enumerate}
    \item $\sum_{j\ge 1} \mathcal{X}_{j} + \mathcal{X}_{0} = 1$,
    \item $\|\mathcal{X}_j\|_\infty +  r_j \|\nabla_x \mathcal{X}_j\|_\infty \le C$ for all $j\ge 1$,
    \item $\textrm{supp}(\mathcal{X}_j) \subset B_n^{j}$ for all $j\ge 0$.
\end{enumerate}
Next, we decompose
\begin{equation*}
    f(x,y) = \sum_{j\ge 1} f(x,y) \mathcal{X}_j(x) + f(x,y)\mathcal{X}_{0}(x):= \sum_{j \ge 1} f_j + f_{0}.
\end{equation*}
Since the covering has finite overlap property (and the finite overlap constant does not depend on $R$), one infers
\begin{align}\label{eqlast}
    \left| \{\mathbb{R}^{n+m}; |f|\ge \lambda\} \right| \le \sum_{j\ge 0} \mathcal{F}_j,
\end{align}
where
\begin{equation*}
    \mathcal{F}_j = \left| \{B_n^{j} \times \mathbb{R}^m; |f_j|\ge C_L^{-1}\lambda\} \right|,\quad \forall j\ge 0.
\end{equation*}


Now, for each $j \ge 1$, we further split (for some $t>0$)
\begin{equation*}
    f_j = \left( f_j - e^{-tL}f_j \right) + e^{-tL}f_j.
\end{equation*}
Note that by \cite[Prop~3.1]{DS2}, we have $\|e^{-tL}\|_{1\to \infty}\lesssim t^{-\mathcal{Q}/2}$. It follows by interpolation that $\|e^{-tL}f_j\|_\infty \le C t^{-\frac{\mathcal{Q}}{2q}} \|f\|_q$. Therefore, by setting $t = \left(2C_L C \|f\|_q/\lambda\right)^{2q/\mathcal{Q}}$, we deduce that
$$
\left|\{|e^{-tL}f_j|\ge (2C_L)^{-1}\lambda\}\right|= 0.
$$
Thus we obtain for $j\ge 1$
\begin{align}\label{eq_thm_main1}
    \mathcal{F}_j \le \left|\left\{B_n^{j} \times \mathbb{R}^m; |f_j - e^{-tL}f_j| \ge (2C_L)^{-1}\lambda  \right\}\right|.
\end{align}

To estimate \eqref{eq_thm_main1}, we write $f_j - e^{tL}f_j = \int_0^t L e^{-sL}f_j ds$. With Minkowski's integral inequality, $\mathcal{F}_j$ is bounded by
\begin{align}\label{eq_thm_main2}
    2C_L\lambda^{-1} \int_0^t \|L e^{-sL} f_j \|_{L^1(B_n^{j} \times \mathbb{R}^m)} ds.
\end{align}
Let $g\in L^\infty$ supported in $B_n^{j} \times \mathbb{R}^m$ with $\|g\|_{\infty} \le 1$. By integration by parts and Hölder's inequality, it is clear that
\begin{align}\label{eq_thm_main3}
    \iint L e^{-sL} f_j(x,y) g(x,y) dx dy &= \int_{\mathbb{R}^m}\int_{B_n^{j}} \nabla_L f_j(x,y) \cdot \nabla_L e^{-sL}g(x,y) dx dy\\ \nonumber
    &\le \left\| \nabla_L f_j \right\|_{L^1(\mathbb{R}^{n+m})} \left\| \nabla_L e^{-tL}g\right\|_{L^\infty(B_n^{j}\times \mathbb{R}^m)}.
\end{align}

\medskip

\textit{Claim.}
\begin{align}\label{claim}
    \sup_{j \ge 1} \sup_{s>0} \|\sqrt{s} \nabla_L e^{-sL}g \|_{L^\infty(B_n^{j}\times \mathbb{R}^m)} \le  C\|g\|_\infty
\end{align}
for all $g\in L^\infty$ supported in $B_n^{j} \times \mathbb{R}^m$, where $C$ depends only on $n,m,\alpha,\beta,C_d,C_D$.

\begin{proof}[Proof of \textit{Claim.}]

Let $\xi = (x,y)\in B_n^j \times\mathbb{R}^m$, $\eta=(x',y') \in B_n^j \times\mathbb{R}^m$ and $s>0$. Note that by Lemma~\ref{le_remoteball}, one has $2^8 r_j \le |x| \le 2^{11} r_j$. Then Lemma~\ref{le_Grushin_gradient} implies that
\begin{align*}
    \left|\nabla_L e^{-sL}g(x,y)\right| &\le \int_{\mathbb{R}^m} \int_{B_n^j} \left(\frac{1}{\sqrt{s}} + \frac{1}{r_j^{1-\alpha}}\right) \frac{C}{V(\xi,\sqrt{s})} e^{-\frac{d(\xi,\eta)^2}{cs}} |g(\eta)| d\eta.
\end{align*}

\textit{Case 1.} If $\sqrt{s}\le 2^{11(1-\alpha)}r_j^{1-\alpha}$, then a standard argument (c.f. \cite[Lemma 2.1]{CD2}) yields that
\begin{align*}
\left|\nabla_L e^{-sL}g(\xi)\right| &\le s^{-1/2} \|g\|_\infty \frac{C}{V(\xi,\sqrt{s})} \int_{\mathbb{R}^{n+m}} e^{-\frac{d(\xi,\eta)^2}{cs}}d\eta \le C s^{-1/2} \|g\|_\infty,
\end{align*}
where the constant $C>0$ depends only on $C_D$.

\textit{Case 2.} If $\sqrt{s}\ge 2^{11(1-\alpha)}r_j^{1-\alpha} \ge |x|^{1-\alpha}$, one deduces
\begin{align}\label{eq_le_gradient1}
\left|\nabla_L e^{-sL}g(\xi)\right| &\le r_j^{\alpha-1} \|g\|_\infty \frac{C}{V(\xi,\sqrt{s})} \int_{B_n^j} \int_{\mathbb{R}^m} e^{-\frac{d(\xi,\eta)^2}{cs}} d y' dx'.
\end{align}
By Lemma~\ref{le_geo}, one has volume estimate $V(\xi,\sqrt{s}) = V((x,y), \sqrt{s}) \sim s^{\mathcal{Q}/2}$. Also, the distance estimate gives
\begin{align*}
    d((x,y);(x',y'))^2 \sim \frac{|x-x'|^2}{\left( |x| + |x'| \right)^{2\alpha}} + \frac{|y-y'|^2}{(|x|+|x'|)^{2\beta} + |y-y'|^{\frac{2\beta}{\beta+1-\alpha}}}.
\end{align*}
Transferring to polar coordinates system (set $\sigma = (|x|+|x'|)^{2\beta}$), the inner integral of \eqref{eq_le_gradient1} is bounded by (up to a constant depending on $C_d$)
\begin{align*}
\int_0^\infty &\textrm{exp}\left(\frac{-cr^2 s^{-1}}{\sigma + r^{2\beta/(\beta+1-\alpha)} } \right) r^{m-1} dr \\
&\le \int_0^{\sigma^{\frac{\beta+1-\alpha}{2\beta}}} \textrm{exp}\left( - \frac{r^2}{2s\sigma} \right) r^{m-1} dr + \int_{\sigma^{\frac{\beta+1-\alpha}{2\beta}}}^\infty \textrm{exp}\left( -\frac{r^{\frac{2-2\alpha}{\beta+1-\alpha}}}{2s} \right) r^{m-1} dr\\
&\le C\max \left(s^{\frac{m}{2}} \sigma^{\frac{m}{2}}, s^{\frac{m(\beta+1-\alpha)}{2(1-\alpha)}} \right) \le C s^{\frac{m(\beta+1-\alpha)}{2(1-\alpha)}},
\end{align*}
where the last inequality follows by observing that $\sigma \sim |x|^{2\beta} \sim r_j^{2\beta} \le C s^{\frac{\beta}{1-\alpha}}$. 

Combining this with \eqref{eq_le_gradient1}, one infers
\begin{align*}
\left|\nabla_L e^{-sL}g(\xi)\right| &\le C\|g\|_\infty r_j^{\alpha-1} s^{-\frac{\mathcal{Q}}{2}} |B_n^j| s^{\frac{m(\beta+1-\alpha)}{2(1-\alpha)}}\\
&\le C \|g\|_\infty r_j^{n+\alpha-1} s^{-\frac{n}{2(1-\alpha)}}\\
&\le Cs^{\frac{n-(1-\alpha)}{2(1-\alpha)} - \frac{n}{2(1-\alpha)}}  \|g\|_\infty  \\
&= C s^{-1/2}\|g\|_\infty
\end{align*}
as desired. This proves the \textit{Claim.}
\end{proof}

\medskip


By the \textit{Claim} \eqref{claim}, \eqref{eq_thm_main2} and \eqref{eq_thm_main1}, one obtains for $j\ge 1$,
\begin{align*}
    \mathcal{F}_j \le C C_L\lambda^{-1} \|\nabla_L f_j\|_1 t^{\frac{1}{2}} &= CC_L \lambda^{-1} \|\nabla_L f_j\|_1 \left(\frac{\|f\|_q}{\lambda}\right)^{\frac{q}{\mathcal{Q}}} \\
    &= CC_L \lambda^{-1-q/\mathcal{Q}} \|\nabla_L f_j\|_1 \|f\|_q^{q/\mathcal{Q}}.
\end{align*}
Note that for $x\in B_n^j$
\begin{align}
    |\nabla_L f_j| &= \left|\left( |x|^\alpha \mathcal{X}_j \nabla_xf  + |x|^\alpha f \nabla \mathcal{X}_j, |x|^\beta \mathcal{X}_j \nabla_yf \right)\right|\\ \nonumber&\le C \left( |\nabla_L f| + |x|^\alpha \frac{|f|}{r_j} \right)\\ \nonumber
    &\le C \left(|\nabla_L f| + \frac{|f|}{|x|^{1-\alpha}}\right).
\end{align}
It then follows by finite overlap property and Lemma~\ref{le_hardy} that
\begin{align*}
\sum_{j \ge 1} \mathcal{F}_j &\le CC_L \lambda^{-1-q/\mathcal{Q}} \|f\|_q^{q/\mathcal{Q}} \left( \sum_{j \ge 1} \|\nabla_L f\|_{L^1(B_n^j \times \mathbb{R}^m)} +  \sum_{j \ge 1} \left\|\frac{f}{|x|^{1-\alpha}}\right\|_{L^1(B_n^j \times \mathbb{R}^m)} \right)\\
&\le CC_L \lambda^{-1-q/\mathcal{Q}} \|f\|_q^{q/\mathcal{Q}} \left( \|\nabla_L f\|_{L^1(\mathbb{R}^{n+m})} + \int_{\mathbb{R}^m} \sum_{j \ge 1} \int_{B_n^j} \frac{|f(x,y)|}{|x|^{1-\alpha}} dx dy \right)\\
&\le CC_L(1+C_H) \lambda^{-1-q/\mathcal{Q}} \|f\|_q^{q/\mathcal{Q}}\|\nabla_L f\|_1.
\end{align*}

To this end, one also needs to treat $\mathcal{F}_0$. By Lemma~\ref{le_hardy}, it is plain that
\begin{align*}
    \mathcal{F}_0 &= \left|\{B_{0}\times \mathbb{R}^m; |f_{0}|\ge C_L^{-1}\lambda\}\right|\\
    &\le C_L \lambda^{-1} \| f \|_{L^1\left(B_n(0,R)\times \mathbb{R}^m\right)} \le C_L \lambda^{-1} R^{1-\alpha} \left\| 
\frac{f}{|x|^{1-\alpha}} \right\|_{L^1(\mathbb{R}^n \times \mathbb{R}^m)}\\
&\le C_LC_H \lambda^{-1-q/\mathcal{Q}} \|f\|_q^{q/\mathcal{Q}}\|\nabla_L f\|_1.
\end{align*}
Consequently, one concludes by \eqref{eqlast} that
\begin{align*}
    \sup_{\lambda>0} \lambda^{1+q/\mathcal{Q}} \left|\left\{ (x,y)\in \mathbb{R}^{n+m}; |f(x,y)|\ge \lambda \right\}\right| \le C \|\nabla_L f\|_1 \|f\|_q^{q/\mathcal{Q}},
\end{align*}
where the constant $C$ only depends on $n,m,\alpha,\beta,C_d,C_L,C_H,C_D$. 

The proof of Theorem~\ref{thm_ISO} now follows from Lemma~\ref{le_BCLS}.
\end{proof}





    









% Therefore, we apply formulas \eqref{carre1}, \eqref{carre2} to $\mathcal{L} = -L$. In what follows, we assume $1\le i,k\le n$, $1\le j,l\le m$.

% Let $f\in C^{\infty}(\mathbb{R}^{n+m})$. Then
% \begin{align*}
%     &\Gamma(f) = \sum_i f_{x_i}^2 + |x|^{2\beta} \sum_j f_{y_j}^2 = |\nabla_L f|^2\\
%     &\mathcal{L}(f)^2 = L(f)^2 = |\Delta_x f|^2 + |x|^{4\beta} |\Delta_y f|^2 + 2|x|^{2\beta} \Delta_x f \Delta_y f,
% \end{align*}
% which are clear.

% Next, we compute $\Gamma_2(f) = \frac{1}{2} \left( \mathcal{L}\Gamma(f) - 2 \Gamma(f,\mathcal{L}f) \right)$.
% \begin{align*}
%     \mathcal{L}\Gamma(f) &= -L\left(\sum_i f_{x_i}^2 + |x|^{2\beta} \sum_j f_{y_j}^2\right)\\ \nonumber
%     &= -\sum_i L\left( f_{x_i}^2 \right) - L \left( |x|^{2\beta} \sum_j f_{y_j}^2 \right)
% \end{align*}
% By Leibuniz's rule, one shows that
% \begin{align}
%     L(gh) = (Lg)h + g(Lh) - 2\nabla_L g \cdot \nabla_L h = (Lg)h + g(Lh) - 2\Gamma(g,h).
% \end{align}
% Therefore,
% \begin{align*}
%     L\left(f_{x_i}^2\right) = 2L \left(f_{x_i} \right) f_{x_i} -2 \Gamma(f_{x_i}),
% \end{align*}
% and
% \begin{align*}
%     L \left( |x|^{2\beta} \sum_j f_{y_j}^2 \right) &= L\left(|x|^{2\beta}\right) \sum_j f_{y_j}^2 + |x|^{2\beta} \sum_j L \left(f_{y_j}^2 \right) -2 \sum_j \Gamma \left(|x|^{2\beta}, f_{y_j}^2 \right)\\ \nonumber
%     &= -2\beta (n+2\beta-2) |x|^{2\beta-2} \sum_j f_{y_j}^2 + 2|x|^{2\beta} \sum_j L(f_{y_j}) f_{y_j}\\ \nonumber
%     &- 2 |x|^{2\beta} \sum_j \Gamma(f_{y_j}) -2 \sum_j \Gamma \left(|x|^{2\beta}, f_{y_j}^2 \right).
% \end{align*}
% Set $S_{xx} = \sum_{i,k=1}^n f_{x_i,x_k}^2$, $S_{xy} = \sum_{i=1}^n \sum_{j=1}^m f_{x_i,y_j}^2$ and $S_{yy} = \sum_{j,l=1}^m f_{y_j,y_l}^2$. One observes that
% \begin{align*}
%     \sum_i \Gamma(f_{x_i}) = \sum_{i,k} f_{x_i,x_k}^2 + |x|^{2\beta} \sum_i\sum_j f_{x_i,y_j}^2 = S_{xx} + |x|^{2\beta} S_{xy},
% \end{align*}
% and
% \begin{align*}
%     \sum_j \Gamma(f_{y_j}) = S_{xy} + |x|^{2\beta} S_{yy}.
% \end{align*}
% One concludes
% \begin{align}\label{appendix1}
%     \mathcal{L}\Gamma(f) &= 2 S_{xx} + 4 |x|^{2\beta} S_{xy} + 2|x|^{4\beta} S_{yy} - 2 \sum_i L(f_{x_i}) f_{x_i} - 2|x|^{2\beta} \sum_j L(f_{y_j}) f_{y_j}\\ \nonumber
%     &+ 2\beta (n+2\beta-2) |x|^{2\beta-2} |\nabla_y f|^2 + 2 \sum_j \Gamma \left(|x|^{2\beta}, f_{y_j}^2 \right).
% \end{align}
% Next, 
% \begin{align*}
%     \Gamma(f,\mathcal{L}f) = -\Gamma(f,Lf) = - \nabla_L f \cdot \nabla_L \left( Lf \right) = - \sum_i f_{x_i} (Lf)_{x_i} - |x|^{2\beta} \sum_j f_{y_j} (Lf)_{y_j}.
% \end{align*}
% Note that
% \begin{align*}
%     (Lf)_{x_i} &= L(f_{x_i}) + 2\beta |x|^{2\beta-2} x_i \Delta_y f.\\
%     (Lf)_{y_j} &= L(f_{y_j}).
% \end{align*}
% We deduce
% \begin{align}
%     \Gamma(f,\mathcal{L}f) = - \sum_i f_{x_i} L(f_{x_i}) - 2\beta |x|^{2\beta-2} \Delta_y f \sum_i x_i f_{x_i} - |x|^{2\beta} \sum_j f_{y_j} L(f_{y_j}).
% \end{align}
% Combine this with \eqref{appendix1} to get
% \begin{align}\label{appendix2}
%     \Gamma_2(f) &= \frac{1}{2} \mathcal{L}\Gamma(f) - \Gamma(f,\mathcal{L}f) \\ \nonumber
%     &= S_{xx} + 2 |x|^{2\beta} S_{xy} + |x|^{4\beta}S_{yy} + \beta (n+2\beta-2) |x|^{2\beta-2} |\nabla_y f|^2 + 2\beta |x|^{2\beta-2} \Delta_y f \sum_i x_i f_{x_i}\\ \nonumber
%     &+ \sum_j \nabla_L \left( |x|^{2\beta} \right) \cdot \nabla_L \left(f_{y_j}^2 \right)\\ \nonumber
%     &= S_{xx} + 2 |x|^{2\beta} S_{xy} + |x|^{4\beta}S_{yy} + \beta (n+2\beta-2) |x|^{2\beta-2} |\nabla_y f|^2 + 2\beta |x|^{2\beta-2} \Delta_y f \sum_i x_i f_{x_i}\\ \nonumber
%     &+ 4\beta |x|^{2\beta-2} \sum_j \sum_i x_i f_{y_j} f_{x_i,y_j}.
% \end{align}




\bigskip
%\newpage
% {\bf Acknowledgments.} 





% {\bf Data availability.}
% This work has neither used nor created new data.


% {\bf Conflict of interest.}
% The author declares that there is no conflict of interest.




\bibliographystyle{abbrv}

\bibliography{references.bib}

% bibliography

\end{document}

